\scanpage{37}
\uline{Motifs sur $\mathbf R$}

$\mathcal M$ \inserted{$\otimes$} catégorie de motifs sur $\mathbf R$, $G=\underline{\operatorname{Aut}}_\otimes(T_B)$ groupe de Galois associé.

Notons qu'on a un $\otimes$-automorphisme d'ordre 2 de $T_B$, qui sur la catégorie $\mathcal V_{\mathbf R}$ des var. proj. lisses $X$ sur $\mathbf R$ revient à l'automorphisme des $H^*(X(\mathbf C),\mathbf Q)$ défini par l'autom $x\mapsto\bar x$ de $X(\mathbf C)$. D'où
\[
\boxed{f_{\deleted{\mathbf R}\infty}\in G(\mathbf Q)}
\qquad(f_{\unsure{\infty|\mathbf R}}^2=1),
\]
(``frobénius réel''). \deleted{On retrouve le sous-espace $H^*_{\mathrm{DR}}(X)$ de $H^*(\unsure{X|X(\mathbf C)},\unsure{\mathbf C|\mathbf Q})=+$} Si $\sigma_M$ est la conjugaison complexe de $M\otimes_{\mathbf Q}\mathbf C$ [où $M\in\operatorname{Ob}\operatorname{Rep}(\unsure{C|G})$], alors la conjugaison $\tau_M$ associée à la structure réelle $T_{\mathrm{DR}}(M)\subset M\otimes_{\mathbf Q}\mathbf C$ n'est autre que
\[
\tau_M=(f_{\unsure{\infty|\mathbf R}})_M\sigma_M
=\sigma_M(f_{\unsure{\infty|\mathbf R}})_M
\]
Donc on retrouve grâce à $f_\infty$ la structure \unsure{mixte|réelle} sur les $M\otimes_{\mathbf R}\mathbf C$, donc aussi le torseur
\[
P=\underline{\operatorname{Isom}}_\otimes(T_B\otimes_{\mathbf Q}\mathbf R,T_{\mathrm{DR}})
\]
sur $G_{\mathbf R}$ : c'est celui défini par le 1-cocycle \deleted{de} \unsure{non abélien} sur $\Gamma=\mathbf Z/2\mathbf Z$, à coeff dans $G(\mathbf C)$, dont la valeur en l'élément $\ne e$ est $f_\infty$. La forme tordue correspondante de $G$ s'obtient en prenant sur $G_{\mathbf C}$ la donnée de descente associée à \deleted{$\unsure{f}$} $\operatorname{int}(f_\infty)\times_{\mathbf R}\sigma_{\mathbf C}$. \deleted{On sait donc, pour} La relation entre $f_\infty$ et
\[
j:S_{\mathbf R}\longrightarrow G_{\mathbf R}
\]
est que
\[
\boxed{f_\infty j(\lambda)f_\infty^{-1}=j(\unsure{\lambda'|\bar\lambda})},
\]
où $\lambda\mapsto\lambda'$ est l'autom can. de $S_{\mathbf R}$ (donnant $\lambda\mapsto\bar\lambda$ sur $S_{\mathbf R}(\mathbf R)\simeq\mathbf C^*$)
\mspage{36}
